\documentclass[10pt,a4paper]{amsart}
\usepackage[margin=19mm]{geometry}
\usepackage{amsmath,amssymb,mathtools}
\usepackage[scr=rsfs,cal=euler]{mathalfa}
\usepackage{xfrac}
\usepackage[unicode,hidelinks,bookmarks=false]{hyperref}
\newcommand{\ie}{i.e.\ }
\newcommand{\cf}{cf.\ }
\newcommand{\resp}{respectively\ }
\newcommand{\Subsection}{\subsection}
\providecommand{\nobreakdash}{\nobreak\hbox{-}}
\setlength{\emergencystretch}{2em}
\multlinegap0pt
\usepackage{bm}
\usepackage{bbm}
\usepackage{dsfont}
\usepackage{xspace}
\usepackage{cancel}
\usepackage{mathtools}
\usepackage{amsthm}
\makeatletter
\@ifundefined{thm}{\newtheorem{thm}{Thm}}{}
\@ifundefined{lem}{\newtheorem{lem}[thm]{Lemma}}{}
\makeatother
\newcommand{\C}{\mathbb{C}}
\newcommand{\G}{\mathbb{G}}
\newcommand{\Q}{\mathbb{Q}}
\newcommand{\Z}{\mathbb{Z}}
\title[A note on polyhedral cones and toric polylogarithms]{A note on polyhedral cones and toric polylogarithms}
\author{Peter Xu}
\date{Source: arXiv:2607.01543v1 (CC BY 4.0)}
\begin{document}
\maketitle

\noindent\emph{Source and licence.} A note on polyhedral cones and toric polylogarithms. Version arXiv:2607.01543v1 (CC BY 4.0), distributed under the Creative Commons Attribution 4.0 International licence, as recorded in the arXiv licence metadata for this exact version and shown on the abstract page https://arxiv.org/abs/2607.01543v1. Full rights notice: licenses/arxiv-2607.01543-CC-BY-4.0.txt.\par\smallskip

\noindent\textit{Editorial scope.} Counted unit: the Lem whose source label is \texttt{None} in the article (source0.tex lines 346--352, source label \texttt{None}), delivered together with the article's own complete original proof of it (source0.tex lines 353--367, one \texttt{proof} environment of 15 lines). No mathematics is written, repaired or re-derived: statement and proof are the article's own text line for line. Required context retained uncounted: none. Every definition, display and label of those lines is retained unchanged. Omitted from this package and not redistributed: 1--345, 368--430. Nothing inside a retained excerpt is omitted or altered: every retained body is delivered exactly as the article's lines are written, so no omission receipt of any kind accompanies this package. Citations: the retained excerpts carry no citation command at all, so no bibliography entry of the article is transcribed and no reference is redistributed. Reference adaptation: none was needed and none was made; every reference inside the delivered unit resolves inside it, so no unresolved reference or citation is produced. Printed slips kept literal: no printed word, symbol, hypothesis or number of the counted statement or of its proof was altered. Layout and class: the article's own class is \texttt{amsart} with packages including inputenc, graphicx, epsfig, bbm, appendix, esvect, tikz-cd, fancyhdr, amsfonts, amsmath, amssymb, mathabx, amsthm, parskip; the wrapper is the lane's pinned \texttt{amsart} editorial wrapper at 10pt on A4, which restates the theorem-like environments the retained text needs and adds the article's own macro definitions that the retained text needs (\texttt{\textbackslash C}, \texttt{\textbackslash G}, \texttt{\textbackslash Q}, \texttt{\textbackslash Z}), with extra wrapper packages (bm, bbm, dsfont, xspace, cancel, mathtools). The delivered numbering is the unit's own. Assets: no retained span contains a figure environment, a graphics-inclusion command, a PDF-page inclusion, a table or a TikZ picture; no image file is copied into this package, the package declares no asset dependency and no image file is redistributed with it. Licence: version arXiv:2607.01543v1 of this article is distributed under the Creative Commons Attribution 4.0 International licence, as recorded in the arXiv licence metadata for this exact version and shown on the abstract page https://arxiv.org/abs/2607.01543v1. A full rights notice is delivered at licenses/arxiv-2607.01543-CC-BY-4.0.txt. Characters: every retained excerpt is pure ASCII, so the delivered engine, XeLaTeX with its default Latin Modern text font, reports no missing character.
    \begin{lem}
        The module $K_i^{(0)}$ in $K(T_\Q)^{(0)}$ is supported only on summands coming from the generic points of \emph{subgroup tori}: that is, ``linear'' embeddings $(\G_m^i)_\Q \to T_\Q$ associated to $i$ by $n$ matrices. In particular, we may identify the term
        \[
            \left(\bigoplus_{x\in T_\Q^{[n-i]}} K_i^M(k(x))\right)^{(0)}= \bigoplus_{A\in } K_i^M(\G_m^{i})^{(0)}
        \]
        in $K(T_\Q)^{(0)}$. 
    \end{lem}

    \begin{proof}    
        Suppose a trace-fixed element $s$ is supported on a union $C$ of dimension-$i$ closed cycles. Then $[a]_*s=s$ is necessarily supported on a subscheme of $[a]_*C$, so we have $C\subset [a]_*C$ for any $a\in \mathbb{N}$. Since $C$ is a union of finitely many irreducible components and each $[a]_*$ is a finite map, we must actually have $C = [a]_*C$ for each $a$, with each of these isogenies permuting the irreducible components. 
        
        We claim the only such cycles are unions of algebraic tori subgroups, as in the lemma statement: in fact, we claim this is even true for the complexification $T_\C$. Let a \emph{torsion subtorus} of $\G_m^n$ to be a subvariety of the form $\zeta + H$ for a subgroup torus $H\subset \G_m^n$ and a torsion point $\zeta$. 
        
        We claim that it suffices to show $C$ is a union of torsion subtori of dimension $i$. Indeed, suppose that we have some torsion subtorus $\zeta+H$ as a component of $C$, for $\zeta$ of minimal order $c$. Then $C$ must contain a component of $[a]^{-1}(\zeta+H)$ for every $a\in \mathbf{N}$, in particular for $a=c^k$ for all $k\ge 0$. These components are precisely $\zeta'+H$ as $\zeta'$ ranges over $c^k$-th roots of $\zeta$; if $c>1$, these cosets are all distinct as $\zeta$ and $k$ vary, because the $[c^k]$-preimages of the class of the $c$-torsion point $\zeta$ in $\G_m^n/H$ are all distinct - this forces $C$ to have infinitely many components, a contradiction. We conclude that $c=1$, and thus $C$ is a union of subgroup tori of dimension $i$.

        We now prove that the components of any $i$-dimensional $C$ are torsion subtori of dimension $i$ for $0\le i \le n$, by induction on $n$; the case $n=0$ is clear. Assuming the inductive hypothesis for $n-1$, we consider a trace-fixed $C$ of dimension $1\le i \le n$; if $i=n$, the result is obvious. Otherwise, if each component of $C$ is contained in a torsion subtorus of dimension $n-1$, the result follows by the inductive hypothesis. 
        
        Finally, if there is a component of $C$ not contained in any torsion subtorus of dimension $n-1$, this means this component contains a point $x=(x_1,\ldots, x_n)$ such that there is no non-zero vector $(v_1,\ldots, v_n)\in \Z^n$ for which $x_1^{v_1}\cdot \ldots \cdot x_n^{v_n}=1$. But then the powers $x^k$ of this point are all also contained in $C$, and these are Zariski dense in $\G_m$: to see this, suppose that
        \[
            f(z_1,\ldots, z_n) = \sum_{\underline{i}\in S\subset \Z_{\ge 0}^n} a_{\underline{i}} z^{\underline{i}}
        \]
        vanishes on all $x^k$. Note that no two terms of the form $x^{k\underline{i}}$ are equal by hypothesis; thus, the $|S|$ by $|S|$ Vandermonde matrix of the first $|S|$ powers of the set $(x^{\underline{i}})_{\underline{i}\in S}$ is invertible, and the linear relations $(f(x^{k \underline{i}}))_{\underline{i}\in S}$ for $k=1,2,\ldots, |S|$ imply that every coefficient $a_{\underline{i}}=0$, meaning that only the zero polynomial vanishes on all powers of $x$. Thus, this case cannot occur.
    \end{proof}

\end{document}
